\subsubsection{The Pieri rules}

Having proved the Littlewood--Richardson rule, let us discuss a few more of
its consequences. As we know from Example \ref{exa.sf.schur-h-e}, the complete
homogeneous symmetric polynomials $h_{n}$ and the elementary symmetric
polynomials $e_{n}$ (for $n\in\left\{  0,1,\ldots,N\right\}  $) are instances
of Schur polynomials. Thus, by appropriately specializing Theorem
\ref{thm.sf.lr-zy}, we can obtain rules for expressing products of the form
$h_{n}s_{\mu}$ and $e_{n}s_{\mu}$ as sums of Schur polynomials. These rules
are known as the \emph{Pieri rules}. To formulate them, we need some more notations:

\begin{definition}
\label{def.sf.strips}Let $\lambda$ and $\mu$ be two $N$-partitions. \medskip

\textbf{(a)} We write $\lambda/\mu$ for the pair $\left(  \mu,\lambda\right)
$. Such a pair is called a \emph{skew partition}. \medskip

\textbf{(b)} We say that $\lambda/\mu$ is a \emph{horizontal strip} if we have
$\mu\subseteq\lambda$ and the Young diagram $Y\left(  \lambda/\mu\right)  $
has no two boxes lying in the same column. \medskip

\textbf{(c)} We say that $\lambda/\mu$ is a \emph{vertical strip} if we have
$\mu\subseteq\lambda$ and the Young diagram $Y\left(  \lambda/\mu\right)  $
has no two boxes lying in the same row. \medskip

Now, let $n\in\mathbb{N}$. \medskip

\textbf{(d)} We say that $\lambda/\mu$ is a \emph{horizontal }$n$\emph{-strip}
if $\lambda/\mu$ is a horizontal strip and satisfies $\left\vert Y\left(
\lambda/\mu\right)  \right\vert =n$. \medskip

\textbf{(e)} We say that $\lambda/\mu$ is a \emph{vertical }$n$\emph{-strip}
if $\lambda/\mu$ is a vertical strip and satisfies $\left\vert Y\left(
\lambda/\mu\right)  \right\vert =n$.
\end{definition}

\Needspace{40pc}

\begin{example}
Let $N=4$. \medskip

\textbf{(a)} If $\lambda=\left(  8,7,4,3\right)  $ and $\mu=\left(
7,4,4,1\right)  $, then we have $\mu\subseteq\lambda$, and the Young diagram
$Y\left(  \lambda/\mu\right)  $ looks as follows:%
\[
Y\left(  \lambda/\mu\right)  =\ydiagram{7+1,4+3,4+0,1+2}\ \ .
\]
From this picture, it is clear that this skew partition $\lambda/\mu$ is a
horizontal strip (and, in fact, a horizontal $6$-strip, since $\left\vert
Y\left(  \lambda/\mu\right)  \right\vert =6$), but not a vertical strip
(since, e.g., there are $3$ boxes in the second row of $Y\left(  \lambda
/\mu\right)  $). \medskip

\textbf{(b)} If $\lambda=\left(  3,3,2,1\right)  $ and $\mu=\left(
2,2,1,0\right)  $, then we have $\mu\subseteq\lambda$, and the Young diagram
$Y\left(  \lambda/\mu\right)  $ looks as follows:%
\[
Y\left(  \lambda/\mu\right)  =\ydiagram{2+1,2+1,1+1,0+1}\ \ .
\]
From this picture, it is clear that this skew partition $\lambda/\mu$ is a
vertical strip (and, in fact, a vertical $4$-strip, since $\left\vert Y\left(
\lambda/\mu\right)  \right\vert =4$), but not a horizontal strip (since there
are $2$ boxes in the third column of $Y\left(  \lambda/\mu\right)  $).
\medskip

\textbf{(c)} If $\lambda=\left(  4,3,1,1\right)  $ and $\mu=\left(
3,2,1,0\right)  $, then we have $\mu\subseteq\lambda$, and the Young diagram
$Y\left(  \lambda/\mu\right)  $ looks as follows:%
\[
Y\left(  \lambda/\mu\right)  =\ydiagram{3+1,2+1,1+0,0+1}\ \ .
\]
From this picture, it is clear that this skew partition $\lambda/\mu$ is both
a horizontal strip (and, in fact, a horizontal $3$-strip) and a vertical strip
(and, in fact, a vertical $3$-strip). \medskip

\textbf{(d)} If $\lambda=\left(  3,3,2,1\right)  $ and $\mu=\left(
1,1,1,1\right)  $, then we have $\mu\subseteq\lambda$, and the Young diagram
$Y\left(  \lambda/\mu\right)  $ looks as follows:%
\[
Y\left(  \lambda/\mu\right)  =\ydiagram{1+2,1+2,1+1,1+0}\ \ .
\]
From this picture, it is clear that this skew partition $\lambda/\mu$ is
neither a horizontal strip nor a vertical strip.
\end{example}

Horizontal and vertical strips can also be characterized in terms of the
entries of the partitions:

\begin{proposition}
\label{prop.sf.strips.entries}Let $\lambda=\left(  \lambda_{1},\lambda
_{2},\ldots,\lambda_{N}\right)  $ and $\mu=\left(  \mu_{1},\mu_{2},\ldots
,\mu_{N}\right)  $ be two $N$-partitions. \medskip

\textbf{(a)} The skew partition $\lambda/\mu$ is a horizontal strip if and
only if we have%
\[
\lambda_{1}\geq\mu_{1}\geq\lambda_{2}\geq\mu_{2}\geq\cdots\geq\lambda_{N}%
\geq\mu_{N}.
\]


\textbf{(b)} The skew partition $\lambda/\mu$ is a vertical strip if and only
if we have%
\[
\mu_{i}\leq\lambda_{i}\leq\mu_{i}+1\ \ \ \ \ \ \ \ \ \ \text{for each }%
i\in\left[  N\right]  .
\]

\end{proposition}

\begin{proof}
See Exercise \ref{exe.sf.strips.entries}.
\end{proof}

We can now state the \emph{Pieri rules}:

\begin{theorem}
[Pieri rules]\label{thm.sf.pieri}Let $n\in\mathbb{N}$. Let $\mu$ be an
$N$-partition. Then: \medskip

\textbf{(a)} We have%
\[
h_{n}s_{\mu}=\sum_{\substack{\lambda\text{ is an }N\text{-partition;}%
\\\lambda/\mu\text{ is a horizontal }n\text{-strip}}}s_{\lambda}.
\]


\textbf{(b)} We have%
\[
e_{n}s_{\mu}=\sum_{\substack{\lambda\text{ is an }N\text{-partition;}%
\\\lambda/\mu\text{ is a vertical }n\text{-strip}}}s_{\lambda}.
\]

\end{theorem}

\begin{example}
Let $N=4$ and $\mu=\left(  2,1,1,0\right)  $. Then: \medskip

\textbf{(a)} Theorem \ref{thm.sf.pieri} \textbf{(a)} (applied to $n=2$) yields%
\[
h_{2}s_{\mu}=\sum_{\substack{\lambda\text{ is an }N\text{-partition;}%
\\\lambda/\mu\text{ is a horizontal }2\text{-strip}}}s_{\lambda}=s_{\left(
2,2,1,1\right)  }+s_{\left(  3,1,1,1\right)  }+s_{\left(  3,2,1,0\right)
}+s_{\left(  4,1,1,0\right)  },
\]
since the $N$-partitions $\lambda$ for which $\lambda/\mu$ is a horizontal
$2$-strip are precisely the four $N$-partitions $\left(  2,2,1,1\right)  $,
$\left(  3,1,1,1\right)  $, $\left(  3,2,1,0\right)  $ and $\left(
4,1,1,0\right)  $. Here are the Young diagrams $Y\left(  \lambda\right)  $ of
these four $N$-partitions $\lambda$ (with the $Y\left(  \mu\right)  $
subdiagram colored red each time):%
\[
\ydiagram{2,2,1,1} *[*(red)]{2,1,1}\ \ ,\ \ \ \ \ \ \ \ \ \ \ydiagram{3,1,1,1}
*[*(red)]{2,1,1}\ \ ,\ \ \ \ \ \ \ \ \ \ \ydiagram{3,2,1,0} *[*(red)]{2,1,1}%
\ \ ,\ \ \ \ \ \ \ \ \ \ \ydiagram{4,1,1,0} *[*(red)]{2,1,1}\ \ .
\]


\textbf{(b)} Theorem \ref{thm.sf.pieri} \textbf{(b)} (applied to $n=2$)
yields
\[
e_{2}s_{\mu}=\sum_{\substack{\lambda\text{ is an }N\text{-partition;}%
\\\lambda/\mu\text{ is a vertical }2\text{-strip}}}s_{\lambda}=s_{\left(
2,2,1,1\right)  }+s_{\left(  2,2,2,0\right)  }+s_{\left(  3,1,1,1\right)
}+s_{\left(  3,2,1,0\right)  },
\]
since the $N$-partitions $\lambda$ for which $\lambda/\mu$ is a vertical
$2$-strip are precisely the four $N$-partitions $\left(  2,2,1,1\right)  $,
$\left(  2,2,2,0\right)  $, $\left(  3,1,1,1\right)  $ and $\left(
3,2,1,0\right)  $. Here are the Young diagrams $Y\left(  \lambda\right)  $ of
these four $N$-partitions $\lambda$ (with the $Y\left(  \mu\right)  $
subdiagram colored red each time):%
\[
\ydiagram{2,2,1,1} *[*(red)]{2,1,1}\ \ ,\ \ \ \ \ \ \ \ \ \ \ydiagram{2,2,2,0}
*[*(red)]{2,1,1}\ \ ,\ \ \ \ \ \ \ \ \ \ \ydiagram{3,1,1,1} *[*(red)]{2,1,1}%
\ \ ,\ \ \ \ \ \ \ \ \ \ \ydiagram{3,2,1,0} *[*(red)]{2,1,1}\ \ .
\]

\end{example}

\begin{proof}
[Proof of Theorem \ref{thm.sf.pieri}.]See Exercise \ref{exe.sf.pieri}.
\end{proof}

\subsubsection{The Jacobi--Trudi identities}

The \emph{Jacobi--Trudi identities} are determinantal formulas expressing a
skew Schur polynomial $s_{\lambda/\mu}$ in terms of complete homogeneous
symmetric polynomials $h_{n}$ or elementary symmetric polynomials $e_{n}$. We
begin with the former:

\begin{theorem}
[First Jacobi--Trudi formula]\label{thm.sf.jt-h}Let $M\in\mathbb{N}$. Let
$\lambda=\left(  \lambda_{1},\lambda_{2},\ldots,\lambda_{M}\right)  $ and
$\mu=\left(  \mu_{1},\mu_{2},\ldots,\mu_{M}\right)  $ be two $M$-partitions
(i.e., weakly decreasing $M$-tuples of nonnegative integers). Then,%
\[
s_{\lambda/\mu}=\det\left(  \left(  h_{\lambda_{i}-\mu_{j}-i+j}\right)
_{1\leq i\leq M,\ 1\leq j\leq M}\right)  .
\]
(Here, $s_{\lambda/\mu}$ is defined as in Definition \ref{def.sf.skew-schur},
where the semistandard tableaux of shape $\lambda/\mu$ are defined as certain
fillings of $Y\left(  \lambda/\mu\right)  :=\left\{  \left(  i,j\right)
\ \mid\ i\in\left[  M\right]  \text{ and }j\in\left[  \lambda_{i}\right]
\setminus\left[  \mu_{i}\right]  \right\}  $, but their entries are still
supposed to be elements of $\left[  N\right]  $. Compare with Remark
\ref{rmk.sf.N-par-vs-M-par}.)
\end{theorem}

\begin{example}
If $M=3$, then Theorem \ref{thm.sf.jt-h} says that%
\[
s_{\lambda/\mu}=\det\left(  \left(  h_{\lambda_{i}-\mu_{j}-i+j}\right)
_{1\leq i\leq3,\ 1\leq j\leq3}\right)  =\det\left(
\begin{array}
[c]{ccc}%
h_{\lambda_{1}-\mu_{1}} & h_{\lambda_{1}-\mu_{2}+1} & h_{\lambda_{1}-\mu
_{3}+2}\\
h_{\lambda_{2}-\mu_{1}-1} & h_{\lambda_{2}-\mu_{2}} & h_{\lambda_{2}-\mu
_{3}+1}\\
h_{\lambda_{3}-\mu_{1}-2} & h_{\lambda_{3}-\mu_{2}-1} & h_{\lambda_{3}-\mu
_{3}}%
\end{array}
\right)  .
\]
For instance, if $\lambda=\left(  4,2,1\right)  $ and $\mu=\left(
1,0,0\right)  $, then%
\begin{align*}
s_{\lambda/\mu}  &  =\det\left(
\begin{array}
[c]{ccc}%
h_{4-1} & h_{4-0+1} & h_{4-0+2}\\
h_{2-1-1} & h_{2-0} & h_{2-0+1}\\
h_{1-1-2} & h_{1-0-1} & h_{1-0}%
\end{array}
\right)  =\det\left(
\begin{array}
[c]{ccc}%
h_{3} & h_{5} & h_{6}\\
h_{0} & h_{2} & h_{3}\\
h_{-2} & h_{0} & h_{1}%
\end{array}
\right) \\
&  =\det\left(
\begin{array}
[c]{ccc}%
h_{3} & h_{5} & h_{6}\\
1 & h_{2} & h_{3}\\
0 & 1 & h_{1}%
\end{array}
\right)  \ \ \ \ \ \ \ \ \ \ \left(  \text{since }h_{0}=1\text{ and }%
h_{-2}=0\right)  .
\end{align*}

\end{example}

The proof of Theorem \ref{thm.sf.jt-h} is not too hard using what we have
learnt about lattice paths in Section \ref{sec.det.comb.lgv}. Here is an
outline (with some details left to exercises):

\begin{proof}
[Proof of Theorem \ref{thm.sf.jt-h} (sketched).]Let us follow Convention
\ref{conv.lgv.digraph}, Definition \ref{def.lgv.lattice} and Definition
\ref{def.lgv.path-tups}. We will work with the digraph $\mathbb{Z}^{2}$. For
each arc $a$ of the digraph $\mathbb{Z}^{2}$, we define an element $w\left(
a\right)  \in\mathcal{P}$ (called the \emph{weight} of $a$) as follows:

\begin{itemize}
\item If $a$ is an east-step $\left(  i,j\right)  \rightarrow\left(
i+1,j\right)  $ with $j\in\left[  N\right]  $, then we set $w\left(  a\right)
:=x_{j}$.

\item If $a$ is any other arc, then we set $w\left(  a\right)  :=1$.
\end{itemize}

For each path $p$ of $\mathbb{Z}^{2}$, define the \emph{weight} $w\left(
p\right)  $ of $p$ by%
\[
w\left(  p\right)  :=\prod_{a\text{ is an arc of }p}w\left(  a\right)  .
\]


Now, it is not hard to see the following (compare with Proposition
\ref{prop.lgv.1-paths.ct}):

\begin{statement}
\textit{Observation 1:} Let $a$ and $c$ be two integers. Then,%
\[
\sum_{\substack{p\text{ is a path}\\\text{from }\left(  a,1\right)  \text{ to
}\left(  c,N\right)  }}w\left(  p\right)  =h_{c-a}.
\]

\end{statement}

See Exercise \ref{exe.sf.jt-h} \textbf{(a)} for a proof of Observation 1.

Next, for each path tuple $\mathbf{p}=\left(  p_{1},p_{2},\ldots,p_{k}\right)
$, let us define the \emph{weight} $w\left(  \mathbf{p}\right)  $ of
$\mathbf{p}$ by%
\[
w\left(  \mathbf{p}\right)  :=w\left(  p_{1}\right)  w\left(  p_{2}\right)
\cdots w\left(  p_{k}\right)  .
\]


Set $k:=M$ (for the sake of convenience). Thus, $\lambda=\left(  \lambda
_{1},\lambda_{2},\ldots,\lambda_{k}\right)  $ and $\mu=\left(  \mu_{1},\mu
_{2},\ldots,\mu_{k}\right)  $.

Define two $k$-vertices $\mathbf{A}=\left(  A_{1},A_{2},\ldots,A_{k}\right)  $
and $\mathbf{B}=\left(  B_{1},B_{2},\ldots,B_{k}\right)  $ by setting%
\[
A_{i}:=\left(  \mu_{i}-i,\ 1\right)  \ \ \ \ \ \ \ \ \ \ \text{and}%
\ \ \ \ \ \ \ \ \ \ B_{i}:=\left(  \lambda_{i}-i,\ N\right)
\ \ \ \ \ \ \ \ \ \ \text{for each }i\in\left[  k\right]  .
\]
It is easy to see that the conditions (\ref{eq.cor.lgv.kpaths.wt-np.xA}),
(\ref{eq.cor.lgv.kpaths.wt-np.yA}), (\ref{eq.cor.lgv.kpaths.wt-np.xB}) and
(\ref{eq.cor.lgv.kpaths.wt-np.yB}) of Corollary \ref{cor.lgv.kpaths.wt-np} are
satisfied. Hence, (\ref{eq.cor.lgv.kpaths.wt-np.claim}) yields%
\begin{align}
&  \det\left(  \left(  \sum_{p:A_{i}\rightarrow B_{j}}w\left(  p\right)
\right)  _{1\leq i\leq k,\ 1\leq j\leq k}\right) \nonumber\\
&  =\sum_{\substack{\mathbf{p}\text{ is a nipat}\\\text{from }\mathbf{A}\text{
to }\mathbf{B}}}w\left(  \mathbf{p}\right)  , \label{pf.thm.sf.jt-h.5}%
\end{align}
where \textquotedblleft$p:A_{i}\rightarrow B_{j}$\textquotedblright\ means
\textquotedblleft$p$ is a path from $A_{i}$ to $B_{j}$\textquotedblright. The
left hand side of this equality is%
\begin{align}
&  \det\left(  \left(  \sum_{p:A_{i}\rightarrow B_{j}}w\left(  p\right)
\right)  _{1\leq i\leq k,\ 1\leq j\leq k}\right) \nonumber\\
&  =\det\left(  \left(  h_{\left(  \lambda_{j}-j\right)  -\left(  \mu
_{i}-i\right)  }\right)  _{1\leq i\leq k,\ 1\leq j\leq k}\right)
\ \ \ \ \ \ \ \ \ \ \left(  \text{by Observation 1}\right) \nonumber\\
&  =\det\left(  \left(  h_{\left(  \lambda_{i}-i\right)  -\left(  \mu
_{j}-j\right)  }\right)  _{1\leq i\leq k,\ 1\leq j\leq k}\right)
\ \ \ \ \ \ \ \ \ \ \left(  \text{by Theorem \ref{thm.det.transp}}\right)
\nonumber\\
&  =\det\left(  \left(  h_{\lambda_{i}-\mu_{j}-i+j}\right)  _{1\leq i\leq
k,\ 1\leq j\leq k}\right) \nonumber\\
&  =\det\left(  \left(  h_{\lambda_{i}-\mu_{j}-i+j}\right)  _{1\leq i\leq
M,\ 1\leq j\leq M}\right)  \label{pf.thm.sf.jt-h.6}%
\end{align}
(since $k=M$). We shall now analyze the right hand side of
(\ref{pf.thm.sf.jt-h.5}). To that purpose, we need to understand the nipats
from $\mathbf{A}$ to $\mathbf{B}$.

We define the \emph{height} of an east-step $\left(  i,j\right)
\rightarrow\left(  i+1,j\right)  $ to be the number $j$. We define the
\emph{height sequence} of a path $p$ to be the sequence of the heights of the
east-steps of $p$ (going from the starting point to the ending point of $p$).
For example, the path shown in Example \ref{exa.lgv.lattice.path.53} has
height sequence $\left(  1,1,1,2,3\right)  $. It is clear that the height
sequence of a path is always weakly increasing.

If $\mathbf{p}=\left(  p_{1},p_{2},\ldots,p_{k}\right)  $ is a nipat from
$\mathbf{A}$ to $\mathbf{B}$, we let $T\left(  \mathbf{p}\right)  $ be the
tableau of shape $Y\left(  \lambda/\mu\right)  $ such that the entries in the
$i$-th row of $T\left(  \mathbf{p}\right)  $ (for each $i\in\left[  k\right]
$) are the entries of the height sequence of $p_{i}$.

\begin{example}
Let $N=6$ and $M=3$ (so that $k=M=3$) and $\lambda=\left(  4,2,1\right)  $ and
$\mu=\left(  1,0,0\right)  $. Here is a nipat $\mathbf{p}$ from $\mathbf{A}$
to $\mathbf{B}$, and the corresponding tableau $T\left(  \mathbf{p}\right)  $:%
\[
\mathbf{p}=%
%TCIMACRO{\TeXButton{tikz nipat}{\raisebox{-6.5pc}{
%% [tikz figure removed]
%}}}%
%BeginExpansion
\raisebox{-6.5pc}{
% [tikz figure removed]
}%
%EndExpansion
\ \ \Longrightarrow\ \ T\left(  \mathbf{p}\right)
=\raisebox{1.4pc}{\ytableaushort{\none 125,25,4}}\ \ .
\]

\end{example}

We could have defined the tableau $T\left(  \mathbf{p}\right)  $ just as
easily for any path tuple $\mathbf{p}$ from $\mathbf{A}$ to $\mathbf{B}$ (not
just for a nipat); however, the case of a nipat is particularly useful,
because it turns out that the tableau $T\left(  \mathbf{p}\right)  $ is
semistandard if and only if $\mathbf{p}$ is a nipat. Moreover, the following
stronger statement holds:

\begin{statement}
\textit{Observation 2:} There is a bijection%
\begin{align*}
\left\{  \text{nipats from }\mathbf{A}\text{ to }\mathbf{B}\right\}   &
\rightarrow\operatorname*{SSYT}\left(  \lambda/\mu\right)  ,\\
\mathbf{p}  &  \mapsto T\left(  \mathbf{p}\right)  .
\end{align*}

\end{statement}

See Exercise \ref{exe.sf.jt-h} \textbf{(b)} for a proof of Observation 2.

It is easy to see that $w\left(  \mathbf{p}\right)  =x_{T\left(
\mathbf{p}\right)  }$ for any nipat $\mathbf{p}$ from $\mathbf{A}$ to
$\mathbf{B}$. Hence,%
\begin{align*}
\sum_{\substack{\mathbf{p}\text{ is a nipat}\\\text{from }\mathbf{A}\text{ to
}\mathbf{B}}}\underbrace{w\left(  \mathbf{p}\right)  }_{=x_{T\left(
\mathbf{p}\right)  }}  &  =\sum_{\substack{\mathbf{p}\text{ is a
nipat}\\\text{from }\mathbf{A}\text{ to }\mathbf{B}}}x_{T\left(
\mathbf{p}\right)  }=\sum_{T\in\operatorname*{SSYT}\left(  \lambda/\mu\right)
}x_{T}\\
&  \ \ \ \ \ \ \ \ \ \ \ \ \ \ \ \ \ \ \ \ \left(
\begin{array}
[c]{c}%
\text{here, we have substituted }T\text{ for }T\left(  \mathbf{p}\right) \\
\text{in the sum, since the map in}\\
\text{Observation 2 is a bijection}%
\end{array}
\right) \\
&  =s_{\lambda/\mu}\ \ \ \ \ \ \ \ \ \ \left(  \text{since }s_{\lambda/\mu
}\text{ is defined to be }\sum_{T\in\operatorname*{SSYT}\left(  \lambda
/\mu\right)  }x_{T}\right)  .
\end{align*}
Thus,%
\begin{align*}
s_{\lambda/\mu}  &  =\sum_{\substack{\mathbf{p}\text{ is a nipat}\\\text{from
}\mathbf{A}\text{ to }\mathbf{B}}}w\left(  \mathbf{p}\right)  =\det\left(
\left(  \sum_{p:A_{i}\rightarrow B_{j}}w\left(  p\right)  \right)  _{1\leq
i\leq k,\ 1\leq j\leq k}\right)  \ \ \ \ \ \ \ \ \ \ \left(  \text{by
(\ref{pf.thm.sf.jt-h.5})}\right) \\
&  =\det\left(  \left(  h_{\lambda_{i}-\mu_{j}-i+j}\right)  _{1\leq i\leq
M,\ 1\leq j\leq M}\right)  \ \ \ \ \ \ \ \ \ \ \left(  \text{by
(\ref{pf.thm.sf.jt-h.6})}\right)  .
\end{align*}
This proves Theorem \ref{thm.sf.jt-h}.
\end{proof}

Our above proof of Theorem \ref{thm.sf.jt-h} is essentially taken from
\cite[First proof of Theorem 7.16.1]{Stanley-EC2}; other proofs can be found
in \cite[Exercise 2.7.13]{GriRei} (see also \cite[paragraph after Theorem
2.4.6]{GriRei} for several references).

The \emph{second Jacobi--Trudi formula} involves elementary symmetric
polynomials $e_{n}$ (instead of $h_{n}$) and transpose partitions (as in
Exercise \ref{exe.pars.transpose}):

\begin{theorem}
[Second Jacobi--Trudi formula]\label{thm.sf.jt-e}Let $\lambda$ and $\mu$ be
two partitions. Let $\lambda^{t}$ and $\mu^{t}$ be the transposes of $\lambda$
and $\mu$. Let $M\in\mathbb{N}$ be such that both $\lambda^{t}$ and $\mu^{t}$
have length $\leq M$. We extend the partitions $\lambda^{t}$ and $\mu^{t}$ to
$M$-tuples (by inserting zeroes at the end). Write these $M$-tuples
$\lambda^{t}$ and $\mu^{t}$ as $\lambda^{t}=\left(  \lambda_{1}^{t}%
,\lambda_{2}^{t},\ldots,\lambda_{M}^{t}\right)  $ and $\mu=\left(  \mu_{1}%
^{t},\mu_{2}^{t},\ldots,\mu_{M}^{t}\right)  $. Then,%
\[
s_{\lambda/\mu}=\det\left(  \left(  e_{\lambda_{i}^{t}-\mu_{j}^{t}%
-i+j}\right)  _{1\leq i\leq M,\ 1\leq j\leq M}\right)  .
\]

\end{theorem}

\begin{example}
If $M=3$, then Theorem \ref{thm.sf.jt-e} says that%
\[
s_{\lambda/\mu}=\det\left(  \left(  e_{\lambda_{i}^{t}-\mu_{j}^{t}%
-i+j}\right)  _{1\leq i\leq M,\ 1\leq j\leq M}\right)  =\det\left(
\begin{array}
[c]{ccc}%
e_{\lambda_{1}^{t}-\mu_{1}^{t}} & e_{\lambda_{1}^{t}-\mu_{2}^{t}+1} &
e_{\lambda_{1}^{t}-\mu_{3}^{t}+2}\\
e_{\lambda_{2}^{t}-\mu_{1}^{t}-1} & e_{\lambda_{2}^{t}-\mu_{2}^{t}} &
e_{\lambda_{2}^{t}-\mu_{3}^{t}+1}\\
e_{\lambda_{3}^{t}-\mu_{1}^{t}-2} & e_{\lambda_{3}^{t}-\mu_{2}^{t}-1} &
e_{\lambda_{3}^{t}-\mu_{3}^{t}}%
\end{array}
\right)  .
\]
For instance, if $\lambda=\left(  3,2,2\right)  $ and $\mu=\left(
1,1,0\right)  $, then $\lambda^{t}=\left(  3,3,1\right)  $ and $\mu
^{t}=\left(  2\right)  =\left(  2,0,0\right)  $ (here, we have extended the
partition $\mu^{t}$ to an $M$-tuple by inserting zeroes at the end), so that
this becomes
\begin{align*}
s_{\lambda/\mu}  &  =\det\left(
\begin{array}
[c]{ccc}%
e_{3-2} & e_{3-0+1} & e_{3-0+2}\\
e_{3-2-1} & e_{3-0} & e_{3-0+1}\\
e_{1-2-2} & e_{1-0-1} & e_{1-0}%
\end{array}
\right)  =\det\left(
\begin{array}
[c]{ccc}%
e_{1} & e_{4} & e_{5}\\
e_{0} & e_{3} & e_{4}\\
e_{-3} & e_{0} & e_{1}%
\end{array}
\right) \\
&  =\det\left(
\begin{array}
[c]{ccc}%
e_{1} & e_{4} & e_{5}\\
1 & e_{3} & e_{4}\\
0 & 1 & e_{1}%
\end{array}
\right)  \ \ \ \ \ \ \ \ \ \ \left(  \text{since }e_{0}=1\text{ and }%
e_{-3}=0\right)  .
\end{align*}

\end{example}

\begin{proof}
[Proof of Theorem \ref{thm.sf.jt-e}.]See Exercise \ref{exe.sf.jt-e}.
\end{proof}

\newpage

\appendix


\section{\label{chap.hw}Homework exercises}

What follows is a collection of problems (of varying difficulty) that are
meant to illuminate, expand upon and otherwise complement the above text.

The numbers in the squares (like \fbox{3}) are the experience points you gain
for solving the problems. They are a mix of difficulty rating and relevance
score: The harder or more important the problem, the larger is the number in
the square. I believe a \fbox{5} represents a good graduate-level homework
problem that requires thinking and work. A \fbox{3} usually requires some
thinking \textbf{or} work. A \fbox{1} is a warm-up question. A \fbox{7} should
be somewhat too hard for regular homework. Anything above \fbox{10} is not
really meant as homework, but I'd be excited to hear your ideas. Multi-part
exercises sometimes have points split between the parts -- i.e., if parts
\textbf{(b)} and \textbf{(c)} of an exercise are solved using the same idea,
then they may both be assigned \fbox{3} points even if each for itself would
be a \fbox{5}.

In solving an exercise, you can freely use (without proof) the claims of all
exercises above it.

Your goal (for an A grade in the 2024 iteration of Math 531) is to gain at
least 20 experience points from each of the Chapters 3--7 (counting Chapter 2
as part of Chapter 3).
